\section{Robustness}\label{sec:robustness}

Six dimensions of robustness feed into the headline: the strict-invariance limiting case; the cost-distribution estimator ($F_c^k$ regime); the kernel bandwidth used in the Convite GPV cross-check; the bid-level sample filter; the temporal window; and the set of alternative policy variants considered. The first five test the identification and equilibrium-selection assumptions; the sixth feeds the policy discussion in Section~\ref{sec:discussion}. The central findings survive them: the within-auction component of the simulated total effect remains the larger of the two components in both classes, with the share remaining in the 73--85 percent range across cost-distribution estimators and across the two specifications of the post-policy SME pool, and the welfare comparison against the 10 percent price preference remains conditional on goods characteristics rather than universal.

\subsection{Strict-primitive-invariance benchmark}

Table~\ref{tab:v3_strict_invariance} reports the limiting case in which $F_c^{\text{SME,Post}}$ is replaced by $F_c^{\text{SME,Pre}}$ while preserving the observed post-policy SME entry count. This is the strict-invariance bound on the equilibrium-selection main specification of Assumption~\ref{a:setaside}. The simulated total price effect is $+0.29$ in non-pharmaceuticals and $+0.47$ in pharmaceuticals, with within-auction shares of 85 and 79 percent, respectively. Welfare loss at $\lambda = 0.30$ is 30.0 percent of $p_{S_1}$ in non-pharmaceuticals and 39.2 percent in pharmaceuticals.

What changes is the policy ranking in the thin, heterogeneous pharma market. In non-pharma, the 10 percent price preference remains Pareto-superior under strict invariance, just as under the main specification. In pharma, the preference ranking reverses and the implied indifference weight (defined formally in equation~\eqref{eq:welfare_weight}, Section~\ref{sec:discussion}) falls to $w^{\text{SME}}_\star = 0.7$, versus 2.61 under the main equilibrium-selection specification. The bifurcation is the diagnostic: it identifies precisely where equilibrium selection among protected bidders is first-order, agreeing across the two readings in thick standardized markets and separating them in thin heterogeneous ones. Read economically, strict invariance imposes that the SMEs entering under the protected regime draw from the same cost distribution as the SMEs that entered under the open regime---the most generous assumption to the policy, since it implies the protective regime selects entrants whose cost primitives match incumbents'. The realistic alternative is that protected entrants draw from the higher-cost tail (firms unprofitable under the open regime that become profitable under exclusion), which is what equilibrium selection in the main specification captures. Strict invariance therefore understates the welfare cost in pharma; the main-specification cost of 47.0 percent of $p_{S_1}$ is the welfare-conservative number relative to the equilibrium-selection mechanism, not an aggressive upper bound.

\subsection{Cost-distribution regime}

The baseline all-bidders estimator biases the upper tail of $F_c^k$ slightly downward; losers-only biases it materially upward; a Turnbull NPMLE that treats the winner as left-censored delivers point identification. Table~\ref{tab:v3_sensitivity_fc} runs the BNE decomposition under each. The three point estimates for $p_{S_3} - p_{S_1}$ are:
\begin{center}
\begin{tabular}{lccc}
\toprule
Class      & Losers-only & All-bidders (baseline) & Turnbull \\
\midrule
Non-pharma & 0.275 & 0.259 & 0.246 \\
Pharma     & 0.347 & 0.308 & 0.357 \\
\bottomrule
\end{tabular}
\end{center}
All cells lie within 16 percent of the all-bidders point estimate (non-pharmaceutical losers-only $+6$ percent; non-pharmaceutical Turnbull $-5$ percent; pharmaceutical losers-only $+13$ percent; pharmaceutical Turnbull $+16$ percent). The bootstrap CIs (Table~\ref{tab:v3_bootstrap_ci}) overlap materially across regimes. Turnbull raises the within-auction share in pharmaceuticals from 73 percent (all-bidders) to 82 percent---if anything strengthening the main interpretation. The all-bidders approximation used throughout the paper is defensible without the Turnbull machinery; readers who prefer point identification should see Turnbull and the same decomposition.

\paragraph{Sensitivity to cost affiliation.} A Gaussian-copula relaxation of the IPV restriction (Online Appendix, Table~OA.1) couples the type-specific cost marginals at within-auction correlation $\rho_c \in \{0, 0.1, 0.2, 0.3\}$. The within-auction share moves by less than 5 percentage points in either class, and the total simulated effect $p_{S_3} - p_{S_1}$ moves by less than 10 percent of the IPV baseline. The 73--85 percent within-auction share range remains valid for $\rho_c \le 0.3$. Higher affiliation levels are inconsistent with the cross-modality consistency check of Table~\ref{tab:v3_cross_modality}, which would reject GPV-vs-drop-out convergence under strong common-values contamination.

\paragraph{Maskin-Riley FPSB calibration of V3.} The 10 percent SME price preference V3 is computed under Vickrey-equivalent translation in the main specification. The Online Appendix (Table~OA.2) reports the equivalent under \citet{maskinriley2000} FPSB equilibrium adjustment: SME bid-function shifts of $-0.05$ and non-SME bid-function shifts of $+0.02$ deliver a net effect on the second-order statistic of $0.22$--$0.27$ percent of $p^{\mathrm{ref}}$, well within Monte Carlo noise of the BNE simulation. The V3 vs.\ V0 policy ranking is preserved in both classes; the conditional welfare ranking is robust to the FPSB adjustment. The entry margin under FPSB is held at the Vickrey-equivalent observed values and is not jointly resolved with the FPSB bidding equilibrium; a fully solved nested entry-bidding equilibrium under FPSB-with-preference is outside the scope of this paper.

\subsection{Kernel bandwidth}

The cross-modality check in Section~\ref{sec:identification} uses GPV inversion on Convite bids, which requires a kernel density estimate $g(b)$. Silverman's rule-of-thumb $h = 1.06 \hat\sigma n^{-1/5}$ is the baseline. Table~\ref{tab:v3_bandwidth_grid} runs the GPV under $h$-factors of 0.5, 0.75, 1.0, 1.5, and 2.0. The resulting range of $c_{0.5}$ (median cost) across the eight (period $\times$ pharma $\times$ type) Convite strata is 1.2--3.6 percent of the baseline in seven strata; the thinnest cell (Convite pharma SME Pre, 1{,}707 auctions) reaches 5.2 percent and is the only cell where bandwidth dispersion exceeds the BNE Monte Carlo noise of 1--2 percent. Bandwidth is not a first-order concern for the identification.

\subsection{Bid-level sample filter}

The baseline filter is $c_\epsilon \le 3$ and $n \ge 2$. Table~\ref{tab:v3_filter_sensitivity} compares four alternatives.
\begin{itemize}
\item \emph{Tight} ($c_\epsilon \le 2$, $n \ge 3$): $p_{S_3} - p_{S_1}$ is 15 percent lower in non-pharma (0.215 vs.\ 0.254) and 8 percent lower in pharma (0.303 vs.\ 0.330).
\item \emph{Very tight} ($c_\epsilon \le 1.5$, $n \ge 3$): 31 and 25 percent lower respectively. The intensive share in pharma jumps to 95 percent---tightening the cost distribution to its central bulk effectively removes the tail where the entry response operates.
\item \emph{Loose} ($c_\epsilon \le 5$): 1 percent higher in non-pharma and 6 percent higher in pharma. The pharma move is concentrated in the extra observations with $c_\epsilon \in (3, 5]$---almost always outliers at the 500--1000 percent-of-reference band, probably data-entry errors at the hundred-times level. The baseline cut is designed to truncate these.
\end{itemize}
The baseline sits between tight and loose in both classes, is inside the bootstrap CI under any reasonable cut, and can be defended directly: the 3x-reference upper bound catches plausible error patterns while preserving 99.6 percent of the data.

\subsection{Temporal window}

Table~\ref{tab:v3_window_sensitivity} varies the symmetric window around the March 2018 cutoff. ``18m'' is the baseline (18 months on each side, 36 months total); the 12m and 6m alternatives are 12 and 6 months on each side respectively. The total effect is $p_{S_3} - p_{S_1} = +0.272$ (non-pharma 18m), $+0.267$ (12m), and $+0.269$ (6m); $+0.292$ (pharma 18m), $+0.310$ (12m), $+0.338$ (6m). The non-pharmaceutical effect is essentially flat across windows; the pharmaceutical effect is monotonically increasing as the window shrinks, consistent with a transitional period in which the within-auction adjustment establishes more quickly than the participation-margin adjustment.

The within-auction share is comparably stable:
\begin{center}
\begin{tabular}{lccc}
\toprule
Class      & 18m    & 12m    & 6m \\
\midrule
Non-pharma & 79.7\% & 75.4\% & 76.4\% \\
Pharma     & 67.0\% & 73.0\% & 73.6\% \\
\bottomrule
\end{tabular}
\end{center}
The 18m baseline sits at the upper end of the non-pharma within-auction range and at the lower end of the pharma range. The within-auction share varies by 4.3 percentage points in non-pharma (75.4--79.7 percent) and 6.6 percentage points in pharma (67.0--73.6 percent) across the three windows, and the within-auction component remains the larger of the two components in all six (class $\times$ window) cells. 18m is the mature, steady-state choice, with shorter windows providing a corroborating short-run check.

\subsection{Alternative policy variants}

The 10 percent price preference in Section~\ref{sec:pareto} is informative partly because the alternative variants reveal what bidder exclusion is buying. Table~\ref{tab:v3_apv} also reports two intermediate variants that serve as policy-design sensitivities:

\begin{itemize}
\item \emph{V1} (partial set-aside): 50 percent of auctions run SME-only at the Post pool, 50 percent open at the Pre pool. $\Delta p$ is 33--36 percent of V0's---roughly a third, reflecting the non-linearity of pool restriction. A partial rule buys roughly proportional fiscal savings but at the cost of a more complex administrative rule.
\item \emph{V2} (entry-isolated counterfactual): full set-aside with the SME pool held at pre-period size (no participation response). $\Delta p$ is 152--157 percent of V0. Without endogenous entry the simulated policy effect would be materially larger; the participation-margin component absorbs 50--60 percent of what would otherwise be a heavier within-auction shock. This is why Section~\ref{sec:results} reads the participation response as a partial offset to a larger latent within-auction shock.
\end{itemize}

Taken together, V1--V3 span the policy design space between the status-quo full set-aside and the open regime. V2 isolates what the entry response is worth; V1 prices a middle-ground design; V3 shows that a well-targeted preference instrument captures the SME-protection value at essentially zero price cost in thick markets and gives a diagnostic benchmark in thin markets.
